\section{实验总结}
这次实验把前面课程中的动作录制、视觉识别、坐标变换、建图导航和移动控制结合在了一起。我们大组分别完成了TonyPi抱物路线、AR直接跟随、地图导航跟随、RGB-D实时跟随和三车队列方案，再通过共用场地和相互调试检查各部分之间的关系。

对第1组而言，难点是抱物后步态容易偏移，单纯重复动作很难保持路线。我们保留抓取结束时的手臂姿态，调整下肢动作，并用标记距离和方位决定前进、纠偏及转弯。地图跟随部分则让我们进一步理解了目标位置、机器人定位和可通行路径并不是同一件事，需要把视觉观测、TF和导航规划衔接起来。

第2组的AR实验说明，标签编号、实际尺寸和相机坐标约定虽然容易被忽略，却直接决定后车能否跟随。第3组的RGB-D方案进一步考虑了目标身份、深度有效性和局部环境；这部分仍需要更充分的实际样本和场景测试。第10组通过三车联调检查了前后车间距、圆弧转弯和顺序停车，同时发现围栏及长时间停留的前车会引发误避障，说明仅有距离阈值还不足以区分跟随对象与其他障碍。

各组的方法有不同的适用条件。AR跟随实现直接，但受可见性影响较大；地图方法能够利用环境规划路径，却依赖地图、定位和传感器覆盖；RGB-D方法减少了对固定标签的依赖，但身份保持和深度质量更难处理；队列方式便于扩展车辆，也会把前车转弯和停车的影响继续传给后车。

我们还需要继续完善抓取与放置确认、目标丢失后的恢复、任务状态联动、矮障碍检测及全队终点停车测试。现有局部试验不能代替一次完整比赛，因此后续应在相同场地重复运行，记录每台车的轨迹、接触、丢失和到达情况，再评价整体效果。

小组合作是本次实验的重要收获。各成员虽然有主要负责的方向，但动作、感知和执行问题往往需要一起判断；各组之间也需要交换前车速度、标签位置和故障现象。通过共同搭建、调试和整理报告，我们更清楚地认识到，多机系统的稳定运行来自各模块之间持续配合，而不只是某一个控制算法。
